\section{Acquisition at the target geometric resolution}\label{sec:resolution}
Let $K\subset\R^q$ be compact, with diameter at most one after fixing the units of the squared task. Put
\[
 \mathcal F_K=\{f:K\to[0,1]:\Lip(f)\le1\},\qquad
 d_K(\nu,\nu')=\sup_{f\in\mathcal F_K}\left|\int f\,d(\nu-\nu')\right|.
\]
The constant one belongs to $\mathcal F_K$, so $d_K$ controls total scored mass as well as location. The measures below are the actual subprobability occupation measures of Section~\ref{sec:geometry}, with physical preparation and unscored calls still charged in the denominator. We require that the predictive feature used on an excursion is an executable Borel function of that excursion and its entering boundary label, with a prescribed predictive initialization at preparation. More generally it suffices that its conditional excursion law has this property. An unrecorded posterior about a nonreset latent quantity may not be inserted into $K$ for free. If such information is required, the boundary state or the common predictive realization must first be supplied and accounted for.

Apply Theorem~\ref{thm:modulus} to the family $(x,f)\in X\times\mathcal F_K$, with score $f(p)$ on a scored call and zero otherwise. Denote the resulting modulus by $K_N^{K}$ and set $\zeta_N=K_N^{K}+B_N(X\times\mathcal F_K)$. Conditional excursion scores lie in $[0,\tau]$, so the theorem applies without changing the physical return envelope. In particular, for every entering label,
\begin{equation}\label{eq:occupation-prefix}
 \sup_x\max_{t\le N}\frac tN
       d_K(\nu_t^x,\nu_\infty^x)\le\zeta_N.
\end{equation}
There is also the converse of Theorem~\ref{thm:modulus}, with the same prefix supremum and constant $\bar\mu+2$. The compactness of $\mathcal F_K$ is needed only for the qualitative compact-family theorem, not for these quantitative inequalities.

\begin{theorem}[Resolution-localized acquired mass]\label{thm:resolution}
Suppose a submeasure $\lambda\le\nu_\infty$ has mass $w>0$ and
\[
 \lambda(B(z,r))\le A(r)
 \quad(z\in\conv K,
       \ 0<r\le r_0).
\]
For any $M\ge1$ and $0<r\le\min(1,r_0)$ with $M A(r)\le w/2$,
\begin{equation}\label{eq:local-mass}
 q_M(\nu_N)\ge\frac w2r^2-2r\,d_K(\nu_N,\nu_\infty)
              \ge\frac w2r^2-2r\zeta_N.
\end{equation}
Thus $\zeta_N\le wr/8$ implies $q_M(\nu_N)\ge wr^2/4$. This conclusion retains the true mass $w$; it does not normalize the submeasure to a probability. It holds for arbitrary concentration profiles $A$, including singular and stratified occupation, not just a prescribed power law.
\end{theorem}
\begin{proof}
Fix any $M$ centers $C\subset\conv K$ and put
$\varphi(p)=\min\{\dist(p,C)^2,r^2\}$. The map $\dist(p,C)$ is one-Lipschitz, and the truncated square is $2r$-Lipschitz. Moreover $0\le\varphi\le r^2$, and $r\le1$ gives $0\le\varphi/(2r)\le r/2\le1$, so $\varphi/(2r)\in\mathcal F_K$. At least $w-MA(r)\ge w/2$ of the submeasure lies at distance at least $r$ from $C$. Therefore
\[
 \int\dist(p,C)^2\,d\nu_N
 \ge\int\varphi\,d\nu_N
 \ge\frac w2r^2-2r\,d_K(\nu_N,\nu_\infty).
\]
Infimize over $C$. Equation~\eqref{eq:occupation-prefix} at $t=N$ proves the remaining assertions. Centers can be restricted to $\conv K$ by Euclidean projection; compactness gives attainment if needed. No density with respect to ambient volume has been assumed.
\end{proof}

\begin{corollary}[Finite-acquisition causal memory matching]\label{cor:resolution-memory}
Consider one fixed bounded squared prediction task and a class of legal observers with at most $M$ retained labels and at most $M$ fixed responses for this scored task (including any scored phase in that interface). Different control labels may share a readout; there are at most $M$ distinct codepoints. Assume the following statements uniformly over every competing observer, under that observer's own executed law: its actual limiting predictive occupation has a submeasure as in Theorem~\ref{thm:resolution} with common $w,A$; its full-history baseline is at least a specified $B_*(N)$; and its test-family prefix modulus is bounded by the declared $\zeta_{N,M}$. Assume also that a fixed $Q$-label exploration attains $B_*(N)$ and has the common predictive realization, global relation-preserving $J$-covers of error $h_J$, report continuity and actual occupation-weighted update stability of Theorem~\ref{thm:geometry}.

For $M\ge2Q$ choose $0<r_M\le\min(1,r_0)$ with $M A(r_M)\le w/2$ and $\zeta_{N,M}\le wr_M/8$. If $h_{\lfloor M/Q\rfloor}\le C_0r_M$, then
\begin{equation}\label{eq:resolution-memory}
 \frac w4r_M^2\le
 \inf_{\gamma:\,|\gamma|\le M}\bigl(\calR_N(\gamma)-B_*(N)\bigr)
 \le\bigl[L_f\sqrt{S_0}(C_0r_M+\eta_{\rm num})+\xi\bigr]^2.
\end{equation}
Here $\eta_{\rm num}$ is the certified per-update numerical or calibration error in the preserved relation, and $\xi$ is the task-readout error. The upper program uses $Q\lfloor M/Q\rfloor\le M$ labels. For the fixed exploration, the checkpoint excess is exactly its acquired $q_M$, and hence has the same order whenever the matching cover and mass conditions hold.
\end{corollary}
\begin{proof}
Condition on the full legal observation history for each competing program. The conditional-mean projection identity splits its risk into its own full-history baseline and the squared error of its finite readout. Randomized readouts cannot reduce that squared error below the best of their at most $M$ state-conditional means. Thus the excess over $B_*(N)$ is at least its own $q_M(\nu_N)$; Theorem~\ref{thm:resolution} applies separately to that law. The counted recursive construction of Theorem~\ref{thm:geometry}, with $J=\lfloor M/Q\rfloor$, supplies the upper bound and all program, clock and initialization conditions. It compares with the same finite-$N$ baseline by assumption; an asymptotic gain is not substituted for that baseline. The checkpoint statement is the same projection identity followed by a nearest-center encoder.
\end{proof}
For $A(r)=Cr^d$, one may take $r_M=(w/(2CM))^{1/d}$ provided $M\ge w/(2C\min(1,r_0)^d)$, as well as the other displayed budget restrictions. If $\zeta_{N,M}\le A_0N^{-\alpha}$ uniformly in the memory range, it suffices that
\[
 N\ge \left(\frac{8A_0}{wr_M}\right)^{1/\alpha},
 \qquad\text{of order }M^{1/(\alpha d)}.
\]
This is a finite-resolution acquisition condition, not an unconditional assertion that every optimal-control problem has that sample complexity. Without a uniform all-competitor witness or an attaining exploration, only the fixed-exploration and individual-observer inequalities hold.

\begin{proposition}[The resolution scale cannot be ignored]\label{prop:resolution-sharp}
Let $\nu$ be uniform measure on $[0,1]$, and let $\nu_M$ put mass $1/M$ on each midpoint $(j-1/2)/M$. Then
\[
 d_K(\nu,\nu_M)\le\frac1{4M},\qquad
 q_M(\nu_M)=0,\qquad q_M(\nu)=\frac1{12M^2}.
\]
Thus a perturbation of order the cell radius can destroy the entire finite-memory lower bound, even in one dimension.
\end{proposition}
\begin{proof}
Couple a uniform point with the midpoint of its length-$1/M$ interval. The expected distance is $1/(4M)$, which bounds every one-Lipschitz test difference. The atomic quantization assertion is immediate. For the uniform measure, a Voronoi interval of length $l$ contributes at least $l^3/12$, attained by its midpoint. Convexity gives $\sum_{j\le M}l_j^3\ge1/M^2$, with equality for equal intervals. The measures also have a direct prepared-experiment interpretation: prepare $V$ uniformly, reveal either $V$ or just its interval label, and perform the squared task against the same hidden $V$. The posterior means have laws $\nu$ and $\nu_M$, respectively. A continuous report is not a free finite-bit representation.
\end{proof}
The upper transport estimate for arbitrary squared distortion would instead lose a constant times $d_K$ independent of the cell radius. The truncation in~\eqref{eq:local-mass} is what retains the factor $r$. Both the small-ball quantization argument and its elementary one-dimensional sharpness calculation are standard in origin~\cite{GrafLuschgy}; the new use here is the marked physical-acquisition certificate at that resolution.
